% Supplied comparison exhibit; author-year references appear in the first column.
\begin{ThreePartTable}
\begin{TableNotes}[flushleft]
\footnotesize
\item \textit{Notes.}
CBT = cognitive behavioural therapy;
ACT = acceptance and commitment therapy;
RCT = randomised controlled trial.
For individual trials, $N$ denotes the randomised sample
unless an outcome-specific baseline sample is indicated.
Li et al.'s review-wide sample includes non-randomised studies;
1,744 is the sample across all meta-analysed RCTs, not the
depression-specific denominator.
Heinz et al.\ measured anxiety using GAD-Q-IV and reported
model-derived Cohen's $d$.
The Headspace outcome is a composite of anxiety, depression,
and stress.
\end{TableNotes}

\begin{longtable}{@{}
    >{\raggedright\arraybackslash}p{0.17\textwidth}
    >{\raggedright\arraybackslash}p{0.12\textwidth}
    >{\raggedright\arraybackslash}p{0.32\textwidth}
    >{\raggedright\arraybackslash}p{\dimexpr0.39\textwidth-6\tabcolsep\relax}
    @{}}

\caption{Characteristics of studies included in the effect-size comparison}
\label{tab:study_comparison}\\
\toprule
\textbf{Study reference} & \textbf{$N$} & \textbf{Description of intervention}
& \textbf{Description of population} \\
\midrule
\endfirsthead

\multicolumn{4}{@{}l}{\textit{Table~\thetable\ (continued)}}\\
\toprule
\textbf{Study reference} & \textbf{$N$} & \textbf{Description of intervention}
& \textbf{Description of population} \\
\midrule
\endhead

\midrule
\multicolumn{4}{@{}l}{\textit{Continued on next page}}\\
\endfoot

\bottomrule
\insertTableNotes
\endlastfoot

Carl et al.\ (2020)
&
256
&
Study used \textbf{Daylight:} automated, self-paced smartphone CBT
covering relaxation, stimulus control, cognitive restructuring,
and imaginal exposure. Daily practice encouraged.
The primary endpoint was at \textbf{six weeks}; the comparison plots use anxiety at \textbf{three weeks} and depression at \textbf{six weeks}. Compared with a waitlist receiving no programme.
&
Adults from the US and UK, aged 18--67, with GAD-7 scores
$\geq 10$ and a positive diagnostic assessment for generalised
anxiety disorder. Mean age: 30.9 years; 68\% female.
Stable medication was permitted; CBT for anxiety during the
preceding year was an exclusion criterion.
\\
\addlinespace

Fitzpatrick et al.\ (2017)
&
70
&
Study used \textbf{Woebot:} automated text-based conversational agent
delivering CBT self-help over \textbf{two weeks}, with up to
20 sessions. Compared with an information-only control
receiving the NIMH ebook \textit{Depression in College Students}.
Follow-up occurred approximately 2--3 weeks after baseline.
&
Young adults aged 18--28 recruited through a university
community social media site who self-identified as having
anxiety and depression symptoms. Mean age: 22.2 years;
67\% female. A formal psychiatric diagnosis was not required.
\\
\addlinespace

Heinz et al.\ (2025)
&
210 overall;
142 depression;
116 anxiety
&
Study used \textbf{Therabot:} text-based generative AI chatbot providing
evidence-based mental health support. Daily prompts during
the first \textbf{four weeks}, followed by four weeks of
optional access without prompts. Compared with a waitlist
without chatbot access. Plotted effects are from four weeks.
&
US adults aged $\geq 18$ screening positive for clinically
significant depression, generalised anxiety, or high risk of
feeding/eating disorders. Mean age: 33.9 years; 59.5\% female.
Comorbidity was permitted; the outcome-specific baseline
samples overlap.
\\
\addlinespace

Li et al.\ (2023)
\newline
Meta-analysis
&
1,744 across 15 RCTs;
depression: $N=921$;
anxiety: $N=826$
&
AI-based conversational agents providing psychotherapy,
psychoeducation, emotional support, or related assistance
through text, voice, or multimodal interfaces. CBT and
integrative approaches were common. Intervention durations
ranged from several minutes to six months across the studies
included in the review; interventions and comparators varied.
&
Heterogeneous populations across multiple countries, including
non-clinical participants, individuals with screened or
self-reported symptoms, and patients with diagnosed mental
or physical conditions. The outcome-specific meta-analyses
included subsets of the 15 RCTs: $N=921$ for depression and
$N=826$ for anxiety.
\\
\addlinespace

Zhong et al.\ (2024)
\newline
Meta-analysis
&
3,477 across
18 RCTs
&
AI chatbot interventions targeting depression and anxiety,
with differing chatbot designs and psychotherapeutic
approaches. \textbf{Treatment duration varied}; the authors
reported the strongest benefits after eight weeks and also
examined three-month follow-up effects.
&
Adult samples from trials assessing depressive and/or
anxiety symptoms. Eligibility and population composition
varied across trials; the pooled sample does not represent
a single clinical population.
\\
\addlinespace

Shreekumar and Vautrey (2024)
&
2,384
&
\textbf{Headspace:} three months of free access to guided
mindfulness meditation, with encouragement to complete a
10-day introductory course. Some treatment arms also
received financial incentives for use during the first two
weeks. Compared with waitlist conditions. The plot uses
the \textbf{two-week} mental-health index.
&
US adults interested in trying meditation, recruited through
Facebook and Instagram. Existing meditators, previous
Headspace users, and participants above the specified income
threshold were excluded. The sample was 86\% female.
Elevated anxiety or depression was not required.
\\


\end{longtable}
\end{ThreePartTable}
